\documentclass[11pt]{article}
\usepackage[a4paper,margin=1in]{geometry}
\usepackage{microtype}
\usepackage{booktabs}
\usepackage{enumitem}
\usepackage{hyperref}
\usepackage{xcolor}
\usepackage{longtable}
\hypersetup{colorlinks=true,linkcolor=black,urlcolor=blue,citecolor=black}
\setlist{nosep}

\newcommand{\ProjectTitle}{Research Proposal Title}
\newcommand{\ScientificStatus}{OPEN / PARTIALLY OPEN / CLOSED / UNCERTAIN}
\newcommand{\ProjectPriority}{PRIORITIZE / CONDITIONAL / DEPRIORITIZE}
\newcommand{\Decision}{ACCEPT / REVISE}
\newcommand{\EvidenceCaveat}{State material access limitations or write ``None known''.}

\title{\ProjectTitle}
\author{Generated with the OhMyTaste research workflow}
\date{\today}

\begin{document}
\maketitle

\noindent\textbf{Scientific Status:} \ScientificStatus\\
\textbf{Project Priority:} \ProjectPriority\\
\textbf{Workflow Decision:} \Decision\\
\textbf{Material evidence caveat:} \EvidenceCaveat

\section{Executive Summary}
Summarize the research problem, proposed mechanism, decisive test, scientific status, project priority, and expected contribution.

\section{Research Background and Scope}
Describe established facts and the local research context. State the demonstrated scope of central results: task family, dataset/system, model class, loss, scale, and evaluation regime when relevant. Distinguish literature facts from synthesis.

\section{Tension Map and Research Gap}
State the unresolved relation, contradiction, dissociation, hidden assumption, boundary inversion, measurement failure, or underdetermined mechanism. Explain why the gap survives the closest prior work.

\section{Closest Related Work}
\begin{longtable}{p{0.17\linewidth}p{0.20\linewidth}p{0.19\linewidth}p{0.17\linewidth}p{0.17\linewidth}}
\toprule
Work & What it establishes & Difference from proposal & Evidence access/status & Collision status \\
\midrule
Closest paper & Established result & Decisive difference & full-text / abstract-only; checked status & direct / strong risk / related \\
\bottomrule
\end{longtable}

\section{Discovery Trace}
Show how the candidate emerged from evidence:
\begin{quote}
Evidence A + Evidence B $\rightarrow$ tension or unexplained relation $\rightarrow$ missing causal distinction $\rightarrow$ hypothesis $\rightarrow$ discriminating prediction $\rightarrow$ minimal decisive experiment.
\end{quote}

\section{Proposed Hypothesis}
State the hypothesis and underlying mechanism. Include a prediction that differs from the strongest alternative explanation.

\section{Proposed Method}
Describe the intervention, model, analysis, or theoretical method at the minimum level needed to test the mechanism.

\section{Experimental Design}
Specify data/system, controls, manipulated variables, measurements, repetitions, and analysis plan. Mark unresolved resource assumptions as \texttt{Unknown}.

\section{Baselines and Metrics}
List closest baselines, tuning requirements, primary metrics, negative controls, and any leakage or fairness checks.

\section{Decisive First Experiment}
State the smallest experiment that can materially change the project decision. Include:
\begin{itemize}
  \item proposed-mechanism prediction;
  \item strongest-alternative prediction;
  \item primary measurement;
  \item explicit rejection or revision trigger.
\end{itemize}

\section{Expected Contributions}
State the scientific claim that would remain after a successful experiment. Avoid relying solely on a metric improvement.

\section{Risks and Alternative Explanations}
Include the strongest prior-art threat, strongest alternative mechanism, feasibility risks, evaluation artifacts, competition/attribution risk, and evidence-access limitations. Keep scientific openness separate from project-strategy concerns.

\section{References}
List complete bibliographic references with DOI and/or arXiv identifiers when available. Do not invent missing metadata.

\end{document}
